\documentclass[10pt,a4paper]{amsart}
\usepackage[margin=19mm]{geometry}
\usepackage{amsmath,amssymb,mathtools}
\usepackage[scr=rsfs,cal=euler]{mathalfa}
\usepackage{xfrac}
\usepackage[unicode,hidelinks,bookmarks=false]{hyperref}
\newcommand{\ie}{i.e.\ }
\newcommand{\cf}{cf.\ }
\newcommand{\resp}{respectively\ }
\newcommand{\Subsection}{\subsection}
\providecommand{\nobreakdash}{\nobreak\hbox{-}}
\setlength{\emergencystretch}{2em}
\multlinegap0pt
\usepackage{bm}
\usepackage{bbm}
\usepackage{dsfont}
\usepackage{xspace}
\usepackage{cancel}
\usepackage{mathtools}
\usepackage{amsthm}
\makeatletter
\@ifundefined{thm}{\newtheorem{thm}{Thm}}{}
\@ifundefined{claim}{\newtheorem{claim}[thm]{Claim}}{}
\@ifundefined{lem}{\newtheorem{lem}[thm]{Lemma}}{}
\@ifundefined{vmatrix}{\newtheorem{vmatrix}{Vmatrix}}{}
\makeatother
\newcommand{\Q}{\mathbb{Q}} % The rational numbers.
\title[Some uniform effective results on Andr\'{e}--Oort for sums o]{Some uniform effective results on Andr\'{e}--Oort for sums of powers in $\mathbb{C}^n$}
\author{Guy Fowler}
\date{Source: arXiv:2405.06456v2 (CC BY 4.0)}
\begin{document}
\maketitle

\noindent\emph{Source and licence.} Some uniform effective results on Andr\'{e}--Oort for sums of powers in $\mathbb{C}^n$. Version arXiv:2405.06456v2 (CC BY 4.0), distributed under the Creative Commons Attribution 4.0 International licence, as recorded in the arXiv licence metadata for this exact version and shown on the abstract page https://arxiv.org/abs/2405.06456v2. Full rights notice: licenses/arxiv-2405.06456-CC-BY-4.0.txt.\par\smallskip

\noindent\textit{Editorial scope.} Counted unit: the Claim whose source label is \texttt{cl10} in the article (source0.tex lines 1093--1095, source label \texttt{cl10}), delivered together with the article's own complete original proof of it (source0.tex lines 1097--1142, one \texttt{proof} environment of 46 lines). No mathematics is written, repaired or re-derived: statement and proof are the article's own text line for line. Required context retained uncounted: lem:smallclass (lines 293--295); eq:det (lines 1071--1078). Every definition, display and label of those lines is retained unchanged. Omitted from this package and not redistributed: 1--292, 296--1070, 1079--1092, 1096--1096, 1143--1478. Nothing inside a retained excerpt is omitted or altered: every retained body is delivered exactly as the article's lines are written, so no omission receipt of any kind accompanies this package. Citations: the retained excerpts carry no citation command at all, so no bibliography entry of the article is transcribed and no reference is redistributed. Reference adaptation: none was needed and none was made; every reference inside the delivered unit resolves inside it, so no unresolved reference or citation is produced. Printed slips kept literal: no printed word, symbol, hypothesis or number of the counted statement or of its proof was altered. Layout and class: the article's own class is \texttt{amsart} with packages including babel, datetime2, amssymb, hyperref, caption, multirow, xcolor; the wrapper is the lane's pinned \texttt{amsart} editorial wrapper at 10pt on A4, which restates the theorem-like environments the retained text needs and adds the article's own macro definitions that the retained text needs (\texttt{\textbackslash Q}), with extra wrapper packages (bm, bbm, dsfont, xspace, cancel, mathtools). The delivered numbering is the unit's own. Assets: no retained span contains a figure environment, a graphics-inclusion command, a PDF-page inclusion, a table or a TikZ picture; no image file is copied into this package, the package declares no asset dependency and no image file is redistributed with it. Licence: version arXiv:2405.06456v2 of this article is distributed under the Creative Commons Attribution 4.0 International licence, as recorded in the arXiv licence metadata for this exact version and shown on the abstract page https://arxiv.org/abs/2405.06456v2. A full rights notice is delivered at licenses/arxiv-2405.06456-CC-BY-4.0.txt. Characters: every retained excerpt is pure ASCII, so the delivered engine, XeLaTeX with its default Latin Modern text font, reports no missing character.
\begin{lem}\label{lem:smallclass}
	Let $x$ be a singular modulus of discriminant $\Delta$. If $\lvert \Delta \rvert < 15$, then $x \in \Q$. If $\lvert \Delta \rvert < 39$, then $h(\Delta)  \leq 3$.
\end{lem}

\begin{align}\label{eq:det}
	\begin{vmatrix}
		1 & 1 & 1 & 1\\
		x_1 & x_2 & x_3 & x_4\\
		y_1 & y_2 & y_3 & y_4\\
		z_1 & z_2 & z_3 & z_4
	\end{vmatrix}= 0.
\end{align}

\begin{claim}\label{cl10}
	If $\Delta_x = 4 \Delta$ and $\Delta_y = \Delta_z = \Delta$, where $\Delta \equiv 1 \bmod 8$, then $h(\Delta) \leq 6$.
\end{claim}

\begin{proof}
Suppose that $\Delta_x = 4 \Delta$ and $\Delta_y = \Delta_z = \Delta$, where $\Delta \equiv 1 \bmod 8$. Assume that $h(\Delta) \geq 7$. In particular, $\Delta \notin \{-7, -15\}$ by Lemma~\ref{lem:smallclass}. Thus there are exactly two subdominant singular moduli of discriminant $\Delta$ and no subdominant singular moduli of discriminant $4 \Delta$. Let $(x_1, y_1, z_1)$ be the unique conjugate of $(x, y, z)$ with $x_1$ dominant. We distinguish two subcases.

\textit{Subcase 1: neither of $y_1, z_1$ is dominant.} In this case, we may take $(x_2, y_2, z_2), (x_3, y_3, z_3)$ to be the unique conjugates such that $y_2, z_3$ are dominant. Then
\begin{align}\label{bd3}
	\lvert x_1 y_2 z_3 \rvert \geq (\exp(2 \pi \lvert \Delta \rvert^{1/2}) - 2079)(\exp(\pi \lvert \Delta \rvert^{1/2}) - 2079)^2.
\end{align}
Since $h(\Delta) \geq 7$, we may take $(x_4, y_4, z_4)$ to be a conjugate with none of $x_4, y_4, z_4$ dominant. Then, for $\sigma \in S_4$ such that $\sigma(2, 3, 4) \neq (1, 2, 3)$, we have that either: $\sigma(2)=1$ and
\begin{align}\label{bd4}
	&\lvert x_{\sigma(2)} y_{\sigma(3)} z_{\sigma(4)} \rvert \nonumber \\ 
	\leq &(\exp(2 \pi \lvert \Delta \rvert^{1/2}) + 2079)(\exp(\pi \lvert \Delta \rvert^{1/2}) + 2079)\nonumber \\ 
	&\left(\exp\left(\frac{\pi \lvert \Delta \rvert^{1/2}}{2}\right) + 2079\right),
\end{align}
or: $\sigma(2) \neq 1$ and
\begin{align}\label{bd5}
	&\lvert x_{\sigma(2)} y_{\sigma(3)} z_{\sigma(4)} \rvert \nonumber \\
	\leq &\left(\exp\left(\frac{2 \pi \lvert \Delta \rvert^{1/2}}{3}\right) + 2079\right)(\exp(\pi \lvert \Delta \rvert^{1/2}) + 2079)^2.
\end{align}
The bounds \eqref{bd3}, \eqref{bd4}, and \eqref{bd5} are incompatible with \eqref{eq:det} if $\lvert \Delta \rvert \geq 10$. By Lemma~\ref{lem:smallclass}, $\lvert \Delta \rvert \leq 9$ is impossible since $z \notin \Q$.

\textit{Subcase 2: one of $y_1, z_1$ is dominant.} Without loss of generality, suppose that $y_1$ is dominant. Let $(x_2, y_2, z_2)$ be the unique conjugate with $z_2$ dominant. There also exists a conjugate $(x_3, y_3, z_3)$ with $y_3$ subdominant and $x_3, z_3$ not dominant. There exists a conjugate $(x_4, y_4, z_4)$ with neither of $y_4, z_4$ dominant or subdominant. Note that $x_4$ is thus not dominant as well. Hence,
\begin{align}\label{bd6}
	&\lvert x_1 y_3 z_2 \rvert \nonumber\\
	\geq &(\exp(2 \pi \lvert \Delta \rvert^{1/2}) - 2079)\left(\exp\left(\frac{\pi \lvert \Delta \rvert^{1/2}}{2}\right) - 2079\right) \nonumber \\
	&(\exp(\pi \lvert \Delta \rvert^{1/2}) - 2079).
\end{align}
Now let $\sigma \in S_4$ be such that $\sigma(2, 3, 4) \neq (1, 3, 2)$. If $\sigma(2) \neq 1$, then
\begin{align}\label{bd7}
	&\lvert x_{\sigma(2)} y_{\sigma(3)} z_{\sigma(4)} \rvert \nonumber \\
	\leq &\left(\exp\left(\frac{2 \pi \lvert \Delta \rvert^{1/2}}{3}\right) + 2079\right)(\exp(\pi \lvert \Delta \rvert^{1/2}) + 2079)^2.
\end{align}
If $\sigma(2) = 1$ and $\sigma(4) = 2$, then
\begin{align}\label{bd8}
	&\lvert x_{\sigma(2)} y_{\sigma(3)} z_{\sigma(4)} \rvert \nonumber \\ 
	\leq &(\exp(2 \pi \lvert \Delta \rvert^{1/2}) + 2079)\left(\exp\left(\frac{\pi \lvert \Delta \rvert^{1/2}}{3}\right) + 2079\right) \nonumber \\
	&(\exp(\pi \lvert \Delta \rvert^{1/2}) + 2079).
\end{align}
If $\sigma(2) = 1$ and $\sigma(4) \neq 2$, then
\begin{align}\label{bd9}
	&\lvert x_{\sigma(2)} y_{\sigma(3)} z_{\sigma(4)} \rvert \nonumber \\
	\leq &(\exp(2 \pi \lvert \Delta \rvert^{1/2}) + 2079)\left(\exp\left(\frac{\pi \lvert \Delta \rvert^{1/2}}{2}\right) + 2079\right)^2.
\end{align}
The bounds \eqref{bd6}, \eqref{bd7}, \eqref{bd8}, and \eqref{bd9} are incompatible with \eqref{eq:det} if $\lvert \Delta \rvert \geq 32$. By Lemma~\ref{lem:smallclass}, $\lvert \Delta \rvert \leq 31$ is impossible since $h(\Delta) \geq 7$ by assumption.

We must therefore have that $h(\Delta) \leq 6$.
\end{proof}

\end{document}
